\section{引言：从 2D 到 3D}

本节课由斯坦福大学助理教授 Jiajun Wu（吴佳俊）主讲，聚焦\textbf{3D 视觉}（3D Vision）。与前面几周讲的卷积神经网络、Transformer、视觉语言模型和生成模型不同，3D 视觉涉及一套全新的表示方式和方法论。

\begin{importantbox}{本课三大主题}
\begin{enumerate}
    \item \textbf{3D 表示}：如何在计算机中表示三维物体的几何形状？有哪些显式/隐式表示方式？
    \item \textbf{3D 深度学习}：如何将深度学习应用于不同的 3D 数据表示？
    \item \textbf{3D 生成与重建}：如何利用 NeRF、3D Gaussian Splatting 等方法从 2D 图像重建或生成 3D 场景？
\end{enumerate}
\end{importantbox}

\begin{figure}[H]
\centering
\includegraphics[width=0.95\textwidth]{slides-images/slide-001.jpg}
\caption{课程标题页：CS231N Lecture 15 --- 3D Vision\protect\footnotemark}
\end{figure}
\footnotetext{Slides 第1页。}

在 2D 图像中，表示方式非常统一——像素矩阵。但 3D 世界中，物体种类繁多、尺度各异，从大型建筑到细微纹理，需要多种不同的表示方式来捕捉几何、纹理、材质等信息。

\begin{figure}[H]
\centering
\includegraphics[width=0.95\textwidth]{slides-images/slide-002.jpg}
\caption{3D 物体的多样性：从城市建筑到精细雕塑，3D 表示需要处理不同尺度和复杂度\protect\footnotemark}
\end{figure}
\footnotetext{Slides 第2页。}

\subsection{为什么 3D 表示很重要}

选择一种 3D 表示方式时，需要考虑多个因素：

\begin{itemize}
    \item \textbf{存储}：像素易于存储为矩阵，但 3D 点云是不规则的，隐式函数如何存储？
    \item \textbf{创建}：如何从图片或语言描述生成新的 3D 形状？
    \item \textbf{编辑操作}：简化、细分、平滑、修复等操作是否方便？
    \item \textbf{渲染}：如何将 3D 物体渲染成 2D 图像？
    \item \textbf{动画}：是否支持人体、动物等的运动模拟？
    \item \textbf{与深度学习的集成}：能否方便地输入神经网络进行训练和推理？
\end{itemize}

\begin{knowledgebox}{3D 视觉的核心问题}
3D 视觉在某种意义上是\textbf{渲染的逆过程}：渲染是从 3D 模型生成 2D 图像，而 3D 视觉是从 2D 图像重建 3D 世界。
\end{knowledgebox}

\subsection{本章小结}

3D 视觉与 2D 图像处理有着本质区别——不存在统一的``像素''式表示。不同的 3D 表示方式各有优劣，选择取决于具体应用场景和与深度学习方法的兼容性。

%% ========== 显式表示 ==========

